\section{Analysis and Discussion}
\subsection{Enhancements}

Within the F-RCS framework, a general assumption is required. Consider the probability distributions produced by $n$ models over an output $y_h$ with unknown safety. After an arbitrary time interval $\Delta(t)$, if $y_h$ is determined to be unsafe, this means that the model state $\lvert \phi_h \rangle$ is unsafe. We require that the set of models $W$ that assign high probability to $y_h$ be trustworthy. Specifically, models in $W$ cannot deny the probability distributions reported before $\Delta(t)$. They also cannot modify the reported distributions. Models in $X$ cannot forge identities to impersonate models in $W$. These conditions can be enforced through cryptographic techniques such as blockchain \citep{Li2023FISCOBCOS, authors2018Fabric}. Such solutions are isomorphic to existing security supervision mechanisms for voting \citep{yang2021priscore}. Concrete designs tailored to F-RCS remain an open direction for future work.

For time cost optimization, we can construct an intermediate algorithm by trading off CS and RCS. One such example is RCS from local coins \citep{zhang2023waterbear}. After $R$ unsuccessful sampling rounds, the algorithm flips a coin to decide whether to return $\bot$ or to enter the trace phase. This design improves efficiency by sacrificing part of the safety. The detailed algorithm is provided in \autoref{app:Local Coin}.

\subsection{Additional Insights on Optimal Safety}

Based on experimental results, we observe a clear pattern. When $n$ is fixed, the SR of RCS often exhibits a maximum for different $\lambda$. For example, \autoref{fig:SR_m} shows a maximum at $\lambda = 8$. Very small values of $\lambda$ and very large values of $\lambda$ both fail to achieve optimal safety. For a fixed $n$, it is unclear both how to determine whether such a maximum exists and how to theoretically identify the value $\lambda_h$ corresponding to it.





